\chapter{Compactness of Riemannian Manifolds and Flows}

\section{Convergence and compactness of manifolds}

\begin{definition}
\label{topping-ch7-pointed-cheeger-gromov-convergence}
\source{topping:page-0080}
\dcref{ch7:1.1}
\group{compactness-manifolds}
A sequence $(\M_i,g_i,p_i)$ of smooth, complete, pointed Riemannian manifolds
converges smoothly to a smooth, complete, pointed Riemannian manifold
$(\M,g,p)$ if there are compact sets $\Omega_i\subset\M$ exhausting $\M$, with
$p\in\operatorname{int}(\Omega_i)$, and smooth maps
\[
\phi_i:\Omega_i\longrightarrow\M_i
\]
that are diffeomorphic onto their images and satisfy $\phi_i(p)=p_i$, such that
\[
\phi_i^*g_i\longrightarrow g.
\]
Here exhaustion means that every compact $K\subset\M$ is contained in
$\Omega_i$ for all sufficiently large $i$, and the convergence is uniform on
each compact $K$ for $\phi_i^*g_i-g$ and all of its covariant derivatives with
respect to any fixed background connection.
\end{definition}

\begin{remark}
\label{topping-ch7-pointed-convergence-example}
\source{topping:page-0081}
\dcref{ch7:1.2}
\group{compactness-manifolds}
If every $(\M_i,g_i)$ is the same pointed manifold $(\mathcal{N},h)$, then the
choice of base points $q$ gives convergence to $(\mathcal{N},h,q)$, whereas a
sequence of base points $s_i$ can give a cylindrical limit.
\end{remark}

\begin{remark}
\label{topping-ch7-compact-limit-possible}
\source{topping:page-0081}
\dcref{ch7:1.2}
\group{compactness-manifolds}
Pointed smooth convergence can have a noncompact limit even when every
manifold $\M_i$ is compact.
\end{remark}

For a smoothly convergent pointed sequence, for every $s>0$ and
$k\in\{0\}\cup\NN$ one has
\[
\sup_{i\in\NN}\sup_{B_{g_i}(p_i,s)}\lvert\nabla^k\Rm(g_i)\rvert<\infty,
\tag{7.1.1}
\]
and
\[
\inf_i\inj(\M_i,g_i,p_i)>0.
\tag{7.1.2}
\]
The notation $\inj(\M_i,g_i,p_i)$ denotes the injectivity radius at $p_i$.

\begin{theorem}
\label{topping-ch7-compactness-manifolds}
\source{topping:page-0081}
\uses{topping-ch7-pointed-cheeger-gromov-convergence, topping-ch7-pointed-convergence-example, topping-ch7-compact-limit-possible}
\dcref{ch7:1.3}
\group{compactness-manifolds}

Let $(\M_i,g_i,p_i)$ be complete, smooth, pointed Riemannian manifolds of a
common dimension $n$. If (7.1.1) and (7.1.2) hold, then some subsequence
converges smoothly in the pointed sense to a complete, smooth, pointed
Riemannian manifold $(\M,g,p)$ of dimension $n$.
\end{theorem}

\begin{remark}
\label{topping-ch7-curvature-control-needed}
\source{topping:page-0082}
\dcref{ch7:1.4}
\group{compactness-manifolds}
Compactness theorems of this type can use only a bound on $\lvert\Rm\rvert$ or
even on $\lvert\Ric\rvert$ when a weaker convergence is sought. In the Ricci
flow applications here, bounds on $\lvert\nabla^k\Rm\rvert$ are available from
the derivative estimates of Theorem~\ref{topping-maxpr-first-derivative-estimate}.
\end{remark}

\begin{remark}
\label{topping-ch7-injectivity-radius-needed}
\source{topping:page-0082}
\dcref{ch7:1.5}
\group{compactness-manifolds}
A uniform positive lower bound for the pointed injectivity radii is necessary;
without it, cylinders can degenerate.
\end{remark}

\begin{remark}
\label{topping-ch7-injectivity-radius-propagation}
\source{topping:page-0082}
\dcref{ch7:1.6}
\group{compactness-manifolds}
Under the curvature bounds in (7.1.1), a positive lower bound for
$\inj(\M_i,g_i,p_i)$ gives a positive lower bound at any $q\in\M_i$ depending
on $d_{g_i}(p_i,q)$ and the curvature bounds.
\end{remark}

\begin{remark}
\label{topping-ch7-compactness-usage}
\source{topping:page-0082}
\dcref{ch7:1.7}
\group{compactness-manifolds}
The applications in this chapter use the stronger global estimates
\[
\sup_i\sup_{x\in\M}\lvert\nabla^k\Rm(g_i)\rvert<\infty
\qquad (k\in\{0\}\cup\NN),
\]
which imply the local bounds required in (7.1.1).
\end{remark}

\section{Convergence and compactness of flows}

\begin{definition}
\label{topping-ch7-flow-convergence}
\source{topping:page-0083}
\dcref{ch7:2.1}
\group{compactness-flows}
Let $-\infty\leq a<0<b\leq\infty$. For smooth families of complete
Riemannian manifolds $(\M_i,g_i(t))$ and $(\M,g(t))$, defined for
$t\in(a,b)$, and points $p_i\in\M_i$, $p\in\M$, define
\[
(\M_i,g_i(t),p_i)\longrightarrow(\M,g(t),p)
\]
if there are compact sets $\Omega_i\subset\M$ exhausting $\M$, with
$p\in\operatorname{int}(\Omega_i)$, and smooth maps
$\phi_i:\Omega_i\to\M_i$ diffeomorphic onto their images and satisfying
$\phi_i(p)=p_i$, such that
\[
\phi_i^*g_i(t)\longrightarrow g(t).
\]
The convergence is uniform on every compact subset of $\M\times(a,b)$ for
the difference and all derivatives, including all time derivatives and all
covariant spatial derivatives with respect to a fixed background connection.
\end{definition}

\begin{remark}
\label{topping-ch7-flow-backward-interval}
\source{topping:page-0083}
\dcref{ch7:2.2}
\group{compactness-flows}
This definition applies on an interval such as $(-\infty,0)$ when the
$i$-th flow is defined on $(-T_i,0)$ and $T_i\to\infty$.
\end{remark}

\begin{theorem}
\label{topping-ch7-compactness-ricci-flows}
\source{topping:page-0083,topping:page-0084}
\dcref{ch7:2.3}
\group{compactness-flows}
Let $\M_i$ be $n$-dimensional manifolds with points $p_i\in\M_i$, and let
$g_i(t)$ be complete Ricci flows on $\M_i$ for
$t\in(a,b)$, where $-\infty\leq a<0<b\leq\infty$. If
\[
\sup_i\sup_{x\in\M_i,\,t\in(a,b)}\lvert\Rm(g_i(t))(x)\rvert<\infty
\]
and
\[
\inf_i\inj(\M_i,g_i(0),p_i)>0,
\]
then a subsequence converges in the pointed flow sense to a complete Ricci
flow $g(t)$ on an $n$-dimensional manifold $\M$, with a base point $p\in\M$.
\end{theorem}

\section{Blowing up at singularities I}

Let $(\M,g(t))$ be a Ricci flow on a closed manifold with maximal interval
$[0,T)$, where $T<\infty$. By Theorem~\ref{topping-ricci-flow-curvature-blows-up},
\[
\sup_{\M}\lvert\Rm\rvert(\cdot,t)\longrightarrow\infty
\qquad (t\uparrow T).
\]
Choose $p_i\in\M$ and $t_i\uparrow T$ so that
\[
\lvert\Rm\rvert(p_i,t_i)
=\sup_{x\in\M,\,t\in[0,t_i]}\lvert\Rm\rvert(x,t),
\]
for example by maximizing $\lvert\Rm\rvert$ on
$\M\times[0,T-1/i]$. Then
$\lvert\Rm\rvert(p_i,t_i)\to\infty$. Define
\[
g_i(t)=\lvert\Rm\rvert(p_i,t_i)\,
g\!\left(t_i+\frac{t}{\lvert\Rm\rvert(p_i,t_i)}\right).
\]
By the scaling law from Section~1.2.3, each $g_i(t)$ is a Ricci flow on
\[
\left[-t_i\lvert\Rm\rvert(p_i,t_i),
(T-t_i)\lvert\Rm\rvert(p_i,t_i)\right).
\]

The normalization and curvature estimates are
\[
\lvert\Rm(g_i(0))\rvert(p_i)=1,
\qquad
\sup_{\M}\lvert\Rm(g_i(t))\rvert\leq 1\quad (t\leq0),
\]
and, by Theorem~\ref{topping-maxpr-curvature-norm-bound}, for $t>0$,
\[
\sup_{\M}\lvert\Rm(g_i(t))\rvert
\leq\frac{1}{1-C(n)t}.
\]
Consequently, for every $a<0$ there is $b=b(n)>0$ such that all $g_i(t)$ are
defined on $(a,b)$ and
\[
\sup_i\sup_{\M\times(a,b)}\lvert\Rm(g_i(t))\rvert<\infty.
\]

If the injectivity-radius estimate
\begin{equation}
\inf_i\inj(\M,g_i(0),p_i)>0
\tag{7.3.1}
\end{equation}
holds, Theorem~\ref{topping-ch7-compactness-ricci-flows} gives, after passing to a subsequence,
\[
(\M,g_i(t),p_i)\longrightarrow(\mathcal{N},\hat g(t),p_\infty)
\]
on $(a,b)$, and in particular produces a singularity-model Ricci flow.
The estimate (7.3.1) is the unresolved input at this stage of the chapter.
